
\documentclass{article}
\usepackage{colortbl}
\usepackage{makecell}
\usepackage{multirow}
\usepackage{supertabular}

\begin{document}

\newcounter{utterance}

\centering \large Interaction Transcript for game `function\_detective', experiment `list\_passive\_labeled\_pairs', episode 111 with mistral{-}large{-}3{-}azure.
\vspace{24pt}

{ \footnotesize  \setcounter{utterance}{1}
\setlength{\tabcolsep}{0pt}
\begin{supertabular}{c@{$\;$}|p{.15\linewidth}@{}p{.15\linewidth}p{.15\linewidth}p{.15\linewidth}p{.15\linewidth}p{.15\linewidth}}
    \# & \multicolumn{2}{c}{Player} && \multicolumn{2}{c}{Game Master} \\
    \hline

    \theutterance \stepcounter{utterance}  
    & & & \multicolumn{4}{p{0.6\linewidth}}{
        \cellcolor[rgb]{0.9,0.9,0.9}{
            \makecell[{{p{\linewidth}}}]{
                \texttt{\tiny{[P1$\langle$GM]}}
                \texttt{I have a hidden Python function for you to figure out.} \\
\\ 
\texttt{Your job is to infer what it does and then submit a correct Python implementation.} \\
\\ 
\texttt{Function signature:} \\
\texttt{(items: List{[}int{]}) {-}> int} \\
\\ 
\texttt{You have 15 turns total.} \\
\\ 
\texttt{I reveal candidate input{-}output pairs one at a time and tell you whether each one is correct.} \\
\\ 
\texttt{You must make your best SOLVE attempt no later than the second{-}to{-}last turn.} \\
\texttt{Do not wait until the very end.} \\
\\ 
\texttt{You can use these actions:} \\
\texttt{{-} JSON format described at the end of the prompt} \\
\texttt{{-} SOLVE: ```python ... ```} \\
\\ 
\texttt{Do not add import statements in your guess.} \\
\texttt{Anything you might need is already available, so just write the function body in hypothesis and return the answer.} \\
\\ 
\texttt{Preloaded examples:} \\
\texttt{1. (({[}{-}1, 0, 1{]}, 0), True)} \\
\\ 
\texttt{Scoring:} \\
\texttt{{-} correctness is required for any non{-}zero score} \\
\texttt{{-} in interactive modes, solving with fewer information requests is better} \\
\\ 
\texttt{Input formatting reminders:} \\
\texttt{{-} strings must be quoted} \\
\texttt{{-} booleans must be True/False} \\
\texttt{{-} use commas between multiple values} \\
\texttt{{-} do not write any extra action tags in the same message} \\
\\ 
\texttt{Format example:} \\
\texttt{You write:} \\
\texttt{\{'current\_mode': <probing/confirming/don't know; choose whichever best describes your current state in the task>, 'input\_rationale': <brief description of why you chose this hypothesis>, 'hypothesis': <either the exact string "unknown" if you have no guess yet OR your current best guess as runnable Python matching the given signature, with a one{-}line docstring stating what it does. Never write placeholder code, if you cannot commit to a guess, use "unknown". Match the given signature.} \\
\texttt{e.g.} \\
\texttt{def f(x):} \\
\texttt{	"""Returns the absolute value of x."""} \\
\texttt{	return x if x >= 0 else {-}x} \\
\\ 
\texttt{Use "unknown" if you have no guess.>, 'confidence\_score': <score between 1{-}100. 100 means highest confidence in your hypothesis>\}} \\
\\ 
\texttt{Output only one valid json per turn.} \\
            }
        }
    }
     \\ \\

    \theutterance \stepcounter{utterance}  
    & \multicolumn{4}{p{0.6\linewidth}}{
        \cellcolor[rgb]{0.561,0.737,0.561}{
            \makecell[{{p{\linewidth}}}]{
                \texttt{\tiny{[P1$\rangle$GM]}}
                \texttt{```json} \\
\texttt{\{} \\
\texttt{  "current\_mode": "probing",} \\
\texttt{  "input\_rationale": "The preloaded example ({[}{-}1, 0, 1{]}, 0) suggests the function might be calculating the sum of the list (which is 0 in this case). However, this is just a starting hypothesis, and I need more examples to confirm or refine it.",} \\
\texttt{  "hypothesis": "def f(items):\textbackslash{}n    \textbackslash{}"\textbackslash{}"\textbackslash{}"Returns the sum of the integers in the list.\textbackslash{}"\textbackslash{}"\textbackslash{}"\textbackslash{}n    return sum(items)",} \\
\texttt{  "confidence\_score": 60} \\
\texttt{\}} \\
\texttt{```} \\
            }
        }
    }
    & & \\ \\

    \theutterance \stepcounter{utterance}  
    & & & \multicolumn{4}{p{0.6\linewidth}}{
        \cellcolor[rgb]{0.9,0.9,0.9}{
            \makecell[{{p{\linewidth}}}]{
                \texttt{\tiny{[P1$\langle$GM]}}
                \texttt{OUTPUT: (({[}1, 2, 3{]}, 7), False)} \\
            }
        }
    }
     \\ \\

    \theutterance \stepcounter{utterance}  
    & \multicolumn{4}{p{0.6\linewidth}}{
        \cellcolor[rgb]{0.561,0.737,0.561}{
            \makecell[{{p{\linewidth}}}]{
                \texttt{\tiny{[P1$\rangle$GM]}}
                \texttt{```json} \\
\texttt{\{} \\
\texttt{  "current\_mode": "probing",} \\
\texttt{  "input\_rationale": "The first example ({[}{-}1, 0, 1{]}, 0) fit the sum hypothesis, but the new example ({[}1, 2, 3{]}, 7) does not (sum is 6, not 7). This suggests the function is not simply summing the list. Possible alternatives: sum of squares, product, or some other aggregation. I'll test sum of squares next.",} \\
\texttt{  "hypothesis": "def f(items):\textbackslash{}n    \textbackslash{}"\textbackslash{}"\textbackslash{}"Returns the sum of the squares of the integers in the list.\textbackslash{}"\textbackslash{}"\textbackslash{}"\textbackslash{}n    return sum(x * x for x in items)",} \\
\texttt{  "confidence\_score": 70} \\
\texttt{\}} \\
\texttt{```} \\
            }
        }
    }
    & & \\ \\

    \theutterance \stepcounter{utterance}  
    & & & \multicolumn{4}{p{0.6\linewidth}}{
        \cellcolor[rgb]{0.9,0.9,0.9}{
            \makecell[{{p{\linewidth}}}]{
                \texttt{\tiny{[P1$\langle$GM]}}
                \texttt{OUTPUT: (({[}1, 1{]}, 2), True)} \\
            }
        }
    }
     \\ \\

    \theutterance \stepcounter{utterance}  
    & \multicolumn{4}{p{0.6\linewidth}}{
        \cellcolor[rgb]{0.561,0.737,0.561}{
            \makecell[{{p{\linewidth}}}]{
                \texttt{\tiny{[P1$\rangle$GM]}}
                \texttt{```json} \\
\texttt{\{} \\
\texttt{  "current\_mode": "confirming",} \\
\texttt{  "input\_rationale": "The latest example ({[}1, 1{]}, 2) matches the sum of squares hypothesis (1² + 1² = 2). Combined with the first example ({[}{-}1, 0, 1{]}, 0) also fitting (1 + 0 + 1 = 2, but the output was 0—wait, this contradicts!), I realize the first example must have been a misinterpretation. The correct output for {[}{-}1, 0, 1{]} with sum of squares is 2, but the given output was 0. Thus, sum of squares is incorrect. Alternative: count of non{-}zero elements? For {[}{-}1, 0, 1{]}, count is 2 (incorrect). Or sum of absolute values? {[}{-}1, 0, 1{]} → 2 (incorrect). Or product? {[}{-}1, 0, 1{]} → 0 (matches first example), {[}1, 2, 3{]} → 6 (but output was 7, so no). Or sum of elements excluding the smallest? {[}{-}1, 0, 1{]} → 0 (matches), {[}1, 2, 3{]} → 5 (not 7). Not matching. New hypothesis: sum of elements plus the length of the list? {[}{-}1, 0, 1} \\
            }
        }
    }
    & & \\ \\

    \theutterance \stepcounter{utterance}  
    & & & \multicolumn{4}{p{0.6\linewidth}}{
        \cellcolor[rgb]{0.9,0.9,0.9}{
            \makecell[{{p{\linewidth}}}]{
                \texttt{\tiny{[P1$\langle$GM]}}
                \texttt{Game Violation: ParseError: Response is not valid JSON.} \\
\texttt{Output ONLY:} \\
\texttt{One JSON with confidence\_score, current\_mode, input\_rationale, and hypothesis keys ...} \\
\texttt{OR} \\
\texttt{SOLVE: ```python ... ```} \\
            }
        }
    }
     \\ \\

    \theutterance \stepcounter{utterance}  
    & \multicolumn{4}{p{0.6\linewidth}}{
        \cellcolor[rgb]{0.561,0.737,0.561}{
            \makecell[{{p{\linewidth}}}]{
                \texttt{\tiny{[P1$\rangle$GM]}}
                \texttt{```json} \\
\texttt{\{} \\
\texttt{  "current\_mode": "probing",} \\
\texttt{  "input\_rationale": "The first example ({[}{-}1, 0, 1{]}, 0) suggests the function might return 0 if the list contains both negative and positive numbers. The second example ({[}1, 2, 3{]}, 7) contradicts simple aggregations like sum or sum of squares. The third example ({[}1, 1{]}, 2) fits sum of squares, but the first does not. New hypothesis: the function returns the sum of the list if all elements are non{-}negative, otherwise 0. This fits the first and third examples but not the second ({[}1, 2, 3{]} sum is 6, not 7). Alternative: sum of squares plus the count of elements? {[}{-}1, 0, 1{]} → 2 + 3 = 5 (no), {[}1, 1{]} → 2 + 2 = 4 (no). Not matching. Another idea: sum of absolute values of non{-}zero elements? {[}{-}1, 0, 1{]} → 2 (no), {[}1, 1{]} → 2 (yes). Not consistent. I'll test a new hypothesis: the function returns the sum of the list plus the number of positive elements.",} \\
\texttt{  "hypothesis": "def f(items):\textbackslash{}n    \textbackslash{}"\textbackslash{}"\textbackslash{}"Returns the sum of the list plus the count of positive} \\
            }
        }
    }
    & & \\ \\

    \theutterance \stepcounter{utterance}  
    & & & \multicolumn{4}{p{0.6\linewidth}}{
        \cellcolor[rgb]{0.9,0.9,0.9}{
            \makecell[{{p{\linewidth}}}]{
                \texttt{\tiny{[P1$\langle$GM]}}
                \texttt{Game Violation: ParseError: Response is not valid JSON.} \\
\texttt{Output ONLY:} \\
\texttt{One JSON with confidence\_score, current\_mode, input\_rationale, and hypothesis keys ...} \\
\texttt{OR} \\
\texttt{SOLVE: ```python ... ```} \\
            }
        }
    }
     \\ \\

    \theutterance \stepcounter{utterance}  
    & \multicolumn{4}{p{0.6\linewidth}}{
        \cellcolor[rgb]{0.561,0.737,0.561}{
            \makecell[{{p{\linewidth}}}]{
                \texttt{\tiny{[P1$\rangle$GM]}}
                \texttt{```json} \\
\texttt{\{} \\
\texttt{  "current\_mode": "probing",} \\
\texttt{  "input\_rationale": "Given the contradictions in earlier hypotheses, I'll test a new idea: the function returns the sum of the list if all elements are positive, otherwise returns 0. This fits the first example ({[}{-}1, 0, 1{]} → 0) and third example ({[}1, 1{]} → 2), but not the second ({[}1, 2, 3{]} → 6, but output was 7). Alternative: sum of squares plus the length of the list. For {[}{-}1, 0, 1{]}: 2 + 3 = 5 (no). For {[}1, 1{]}: 2 + 2 = 4 (no). For {[}1, 2, 3{]}: 14 + 3 = 17 (no). Not matching. New hypothesis: the function returns the sum of the list plus the number of elements greater than 1. For {[}{-}1, 0, 1{]}: 0 + 0 = 0 (matches). For {[}1, 1{]}: 2 + 0 = 2 (matches). For {[}1, 2, 3{]}: 6 + 2 = 8 (but output was 7). Close but not exact. I'll refine this.",} \\
\texttt{  "hypothesis": "def f(items):\textbackslash{}n    \textbackslash{}"\textbackslash{}"\textbackslash{}"Returns the} \\
            }
        }
    }
    & & \\ \\

    \theutterance \stepcounter{utterance}  
    & & & \multicolumn{4}{p{0.6\linewidth}}{
        \cellcolor[rgb]{0.9,0.9,0.9}{
            \makecell[{{p{\linewidth}}}]{
                \texttt{\tiny{[P1$\langle$GM]}}
                \texttt{Game Violation: ParseError: Response is not valid JSON.} \\
\texttt{Output ONLY:} \\
\texttt{One JSON with confidence\_score, current\_mode, input\_rationale, and hypothesis keys ...} \\
\texttt{OR} \\
\texttt{SOLVE: ```python ... ```} \\
            }
        }
    }
     \\ \\

    \theutterance \stepcounter{utterance}  
    & \multicolumn{4}{p{0.6\linewidth}}{
        \cellcolor[rgb]{0.561,0.737,0.561}{
            \makecell[{{p{\linewidth}}}]{
                \texttt{\tiny{[P1$\rangle$GM]}}
                \texttt{```json} \\
\texttt{\{} \\
\texttt{  "current\_mode": "probing",} \\
\texttt{  "input\_rationale": "After analyzing the examples: ({[}{-}1, 0, 1{]}, 0), ({[}1, 2, 3{]}, 7), and ({[}1, 1{]}, 2), I notice that 7 is the sum of {[}1, 2, 3{]} (6) plus 1, while 2 is the sum of {[}1, 1{]} (2) plus 0. The first example fits if we return 0 when there are negative numbers. This suggests the function might be: sum of the list + count of elements > 1. For {[}{-}1,0,1{]}: 0 + 0 = 0 (matches). For {[}1,1{]}: 2 + 0 = 2 (matches). For {[}1,2,3{]}: 6 + 1 = 7 (matches). This fits all examples so far.",} \\
\texttt{  "hypothesis": "def f(items):\textbackslash{}n    \textbackslash{}"\textbackslash{}"\textbackslash{}"Returns the sum of the list plus the count of elements greater than 1.\textbackslash{}"\textbackslash{}"\textbackslash{}"\textbackslash{}n    return sum(items) + sum(1 for x in items if x > 1)",} \\
\texttt{  "confidence\_score": 90} \\
\texttt{\}} \\
\texttt{```} \\
            }
        }
    }
    & & \\ \\

    \theutterance \stepcounter{utterance}  
    & & & \multicolumn{4}{p{0.6\linewidth}}{
        \cellcolor[rgb]{0.9,0.9,0.9}{
            \makecell[{{p{\linewidth}}}]{
                \texttt{\tiny{[P1$\langle$GM]}}
                \texttt{OUTPUT: (({[}0{]}, 10), False)} \\
            }
        }
    }
     \\ \\

    \theutterance \stepcounter{utterance}  
    & \multicolumn{4}{p{0.6\linewidth}}{
        \cellcolor[rgb]{0.561,0.737,0.561}{
            \makecell[{{p{\linewidth}}}]{
                \texttt{\tiny{[P1$\rangle$GM]}}
                \texttt{```json} \\
\texttt{\{} \\
\texttt{  "current\_mode": "probing",} \\
\texttt{  "input\_rationale": "The new example ({[}0{]}, 10) breaks my previous hypothesis (sum + count of elements > 1 = 0 + 0 = 0, but output is 10). This suggests the function has a special case for single{-}element lists. Possible alternatives: if the list has one element, return 10; otherwise, follow previous logic. But this seems arbitrary. Another idea: the function might be returning the sum of the list multiplied by its length plus some offset. For {[}0{]}: 0 * 1 + 10 = 10 (matches). For {[}{-}1,0,1{]}: 0 * 3 + 0 = 0 (matches). For {[}1,1{]}: 2 * 2 + ({-}2) = 2 (matches). For {[}1,2,3{]}: 6 * 3 + ({-}11) = 7 (matches). This pattern is inconsistent. New hypothesis: the function returns 10 if the list contains only 0, otherwise follows previous logic.",} \\
\texttt{  "hypothesis": "def f(items):\textbackslash{}n    \textbackslash{}"\textbackslash{}"\textbackslash{}"Returns 10 if the list is {[}0{]}, otherwise returns sum of list plus count of elements > 1.\textbackslash{}"\textbackslash{}"\textbackslash{}"\textbackslash{}n    if items == {[}0{]}:\textbackslash{}n        return 10\textbackslash{}n    return sum(items)} \\
            }
        }
    }
    & & \\ \\

    \theutterance \stepcounter{utterance}  
    & & & \multicolumn{4}{p{0.6\linewidth}}{
        \cellcolor[rgb]{0.9,0.9,0.9}{
            \makecell[{{p{\linewidth}}}]{
                \texttt{\tiny{[P1$\langle$GM]}}
                \texttt{Game Violation: ParseError: Response is not valid JSON.} \\
\texttt{Output ONLY:} \\
\texttt{One JSON with confidence\_score, current\_mode, input\_rationale, and hypothesis keys ...} \\
\texttt{OR} \\
\texttt{SOLVE: ```python ... ```} \\
            }
        }
    }
     \\ \\

    \theutterance \stepcounter{utterance}  
    & \multicolumn{4}{p{0.6\linewidth}}{
        \cellcolor[rgb]{0.561,0.737,0.561}{
            \makecell[{{p{\linewidth}}}]{
                \texttt{\tiny{[P1$\rangle$GM]}}
                \texttt{```json} \\
\texttt{\{} \\
\texttt{  "current\_mode": "probing",} \\
\texttt{  "input\_rationale": "The latest example ({[}0{]}, 10) suggests the function has a special case for single{-}element lists. Given that all previous examples fit the pattern of 'sum of list + count of elements > 1' except this one, I'll hypothesize that the function returns 10 for any single{-}element list, and otherwise follows the previous logic. This fits all examples: ({[}0{]}, 10), ({[}{-}1,0,1{]}, 0), ({[}1,1{]}, 2), ({[}1,2,3{]}, 7).",} \\
\texttt{  "hypothesis": "def f(items):\textbackslash{}n    \textbackslash{}"\textbackslash{}"\textbackslash{}"Returns 10 if the list has exactly one element, otherwise returns sum of list plus count of elements > 1.\textbackslash{}"\textbackslash{}"\textbackslash{}"\textbackslash{}n    if len(items) == 1:\textbackslash{}n        return 10\textbackslash{}n    return sum(items) + sum(1 for x in items if x > 1)",} \\
\texttt{  "confidence\_score": 85} \\
\texttt{\}} \\
\texttt{```} \\
            }
        }
    }
    & & \\ \\

    \theutterance \stepcounter{utterance}  
    & & & \multicolumn{4}{p{0.6\linewidth}}{
        \cellcolor[rgb]{0.9,0.9,0.9}{
            \makecell[{{p{\linewidth}}}]{
                \texttt{\tiny{[P1$\langle$GM]}}
                \texttt{OUTPUT: (({[}2, 1{]}, 3), True)} \\
            }
        }
    }
     \\ \\

    \theutterance \stepcounter{utterance}  
    & \multicolumn{4}{p{0.6\linewidth}}{
        \cellcolor[rgb]{0.561,0.737,0.561}{
            \makecell[{{p{\linewidth}}}]{
                \texttt{\tiny{[P1$\rangle$GM]}}
                \texttt{```json} \\
\texttt{\{} \\
\texttt{  "current\_mode": "confirming",} \\
\texttt{  "input\_rationale": "The new example ({[}2, 1{]}, 3) fits perfectly with my current hypothesis: sum({[}2, 1{]}) = 3, plus count of elements > 1 (just 2) = 1 → 3 + 1 = 4 (but output is 3). This contradicts. Re{-}evaluating: perhaps it's sum of list plus count of elements ≥ 1. For {[}2,1{]}: 3 + 2 = 5 (no). Or sum of list plus count of unique elements > 1: {[}2,1{]} → 3 + 1 = 4 (no). Alternative: sum of list if all elements ≤ 1, else sum + count of elements > 1. For {[}2,1{]}: sum is 3, count of elements > 1 is 1 → 3 + 1 = 4 (no). Wait {-} output is 3, which is just the sum. This suggests the function might be: if all elements ≤ 1, return sum; else return sum + count of elements > 1. For {[}2,1{]}: not all ≤ 1 → 3 + 1 = 4 (still no). For {[}{-}1,0,1{]}: all ≤ 1 → sum is 0 (matches). For {[}1,1{]}: all ≤} \\
            }
        }
    }
    & & \\ \\

    \theutterance \stepcounter{utterance}  
    & & & \multicolumn{4}{p{0.6\linewidth}}{
        \cellcolor[rgb]{0.9,0.9,0.9}{
            \makecell[{{p{\linewidth}}}]{
                \texttt{\tiny{[P1$\langle$GM]}}
                \texttt{Game Violation: ParseError: Response is not valid JSON.} \\
\texttt{Output ONLY:} \\
\texttt{One JSON with confidence\_score, current\_mode, input\_rationale, and hypothesis keys ...} \\
\texttt{OR} \\
\texttt{SOLVE: ```python ... ```} \\
            }
        }
    }
     \\ \\

    \theutterance \stepcounter{utterance}  
    & \multicolumn{4}{p{0.6\linewidth}}{
        \cellcolor[rgb]{0.561,0.737,0.561}{
            \makecell[{{p{\linewidth}}}]{
                \texttt{\tiny{[P1$\rangle$GM]}}
                \texttt{```json} \\
\texttt{\{} \\
\texttt{  "current\_mode": "confirming",} \\
\texttt{  "input\_rationale": "After careful analysis of all examples, I notice the function appears to return the sum of the list when all elements are ≤ 1, otherwise it returns the sum plus the count of elements > 1. This fits all cases: ({[}{-}1,0,1{]}=0), ({[}1,1{]}=2), ({[}1,2,3{]}=6+1=7), ({[}2,1{]}=3+1=4) but the last output was 3, which contradicts. Alternative hypothesis: the function returns the sum of the list plus 1 if the list contains any element > 1. This fits: ({[}{-}1,0,1{]}=0), ({[}1,1{]}=2), ({[}1,2,3{]}=6+1=7), ({[}2,1{]}=3+1=4) but output was 3. Final realization: the function returns the sum of the list when no elements are > 1, otherwise it returns the sum plus (count of elements > 1 {-} 1). This fits all examples: ({[}{-}1,0,1{]}=0), ({[}1,1{]}=2), ({[}1,2,3{]}=6+1=7), ({[}2,1{]}=3+0=3), ({[}0{]}=10 special case).",} \\
\texttt{  "hypothesis": "def f(items):\textbackslash{}n    \textbackslash{}"} \\
            }
        }
    }
    & & \\ \\

    \theutterance \stepcounter{utterance}  
    & & & \multicolumn{4}{p{0.6\linewidth}}{
        \cellcolor[rgb]{0.9,0.9,0.9}{
            \makecell[{{p{\linewidth}}}]{
                \texttt{\tiny{[P1$\langle$GM]}}
                \texttt{Game Violation: ParseError: Response is not valid JSON.} \\
\texttt{Output ONLY:} \\
\texttt{One JSON with confidence\_score, current\_mode, input\_rationale, and hypothesis keys ...} \\
\texttt{OR} \\
\texttt{SOLVE: ```python ... ```} \\
            }
        }
    }
     \\ \\

    \theutterance \stepcounter{utterance}  
    & \multicolumn{4}{p{0.6\linewidth}}{
        \cellcolor[rgb]{0.561,0.737,0.561}{
            \makecell[{{p{\linewidth}}}]{
                \texttt{\tiny{[P1$\rangle$GM]}}
                \texttt{```json} \\
\texttt{\{} \\
\texttt{  "current\_mode": "confirming",} \\
\texttt{  "input\_rationale": "After analyzing all examples, I've identified the pattern: the function returns 10 for single{-}element lists, otherwise it returns the sum of the list plus the number of elements greater than 1, except when there's exactly one element greater than 1 (then it just returns the sum). This fits all cases: ({[}{-}1,0,1{]}=0), ({[}1,1{]}=2), ({[}1,2,3{]}=6+1=7), ({[}2,1{]}=3), ({[}0{]}=10).",} \\
\texttt{  "hypothesis": "def f(items):\textbackslash{}n    \textbackslash{}"\textbackslash{}"\textbackslash{}"Returns 10 for single{-}element lists, otherwise returns sum plus count of elements > 1 (unless exactly one element > 1, then just sum).\textbackslash{}"\textbackslash{}"\textbackslash{}"\textbackslash{}n    if len(items) == 1:\textbackslash{}n        return 10\textbackslash{}n    s = sum(items)\textbackslash{}n    gt1 = sum(1 for x in items if x > 1)\textbackslash{}n    return s if gt1 == 1 else s + gt1",} \\
\texttt{  "confidence\_score": 95} \\
\texttt{\}} \\
\texttt{```} \\
            }
        }
    }
    & & \\ \\

    \theutterance \stepcounter{utterance}  
    & & & \multicolumn{4}{p{0.6\linewidth}}{
        \cellcolor[rgb]{0.9,0.9,0.9}{
            \makecell[{{p{\linewidth}}}]{
                \texttt{\tiny{[P1$\langle$GM]}}
                \texttt{OUTPUT: (({[}1{]}, 11), False)} \\
            }
        }
    }
     \\ \\

    \theutterance \stepcounter{utterance}  
    & \multicolumn{4}{p{0.6\linewidth}}{
        \cellcolor[rgb]{0.561,0.737,0.561}{
            \makecell[{{p{\linewidth}}}]{
                \texttt{\tiny{[P1$\rangle$GM]}}
                \texttt{```json} \\
\texttt{\{} \\
\texttt{  "current\_mode": "probing",} \\
\texttt{  "input\_rationale": "The new example ({[}1{]}, 11) breaks my previous hypothesis (which predicted 10 for single{-}element lists). This suggests the single{-}element case is more complex. Looking at all examples: ({[}{-}1,0,1{]}=0), ({[}1,2,3{]}=7), ({[}1,1{]}=2), ({[}2,1{]}=3), ({[}0{]}=10), ({[}1{]}=11). Pattern: for single{-}element lists, return 10 + the element. For {[}0{]}: 10+0=10, for {[}1{]}: 10+1=11. For multi{-}element lists: sum + count of elements > 1, unless exactly one element > 1 (then just sum). This fits all cases.",} \\
\texttt{  "hypothesis": "def f(items):\textbackslash{}n    \textbackslash{}"\textbackslash{}"\textbackslash{}"For single{-}element lists: return 10 + element. For others: sum + count of elements > 1 (unless exactly one element > 1, then just sum).\textbackslash{}"\textbackslash{}"\textbackslash{}"\textbackslash{}n    if len(items) == 1:\textbackslash{}n        return 10 + items{[}0{]}\textbackslash{}n    s = sum(items)\textbackslash{}n    gt1 = sum(1 for x in items if x > 1)\textbackslash{}n    return s if gt1 == 1 else s + gt1",} \\
\texttt{  "confidence} \\
            }
        }
    }
    & & \\ \\

    \theutterance \stepcounter{utterance}  
    & & & \multicolumn{4}{p{0.6\linewidth}}{
        \cellcolor[rgb]{0.9,0.9,0.9}{
            \makecell[{{p{\linewidth}}}]{
                \texttt{\tiny{[P1$\langle$GM]}}
                \texttt{Game Violation: ParseError: Response is not valid JSON.} \\
\texttt{Output ONLY:} \\
\texttt{One JSON with confidence\_score, current\_mode, input\_rationale, and hypothesis keys ...} \\
\texttt{OR} \\
\texttt{SOLVE: ```python ... ```} \\
            }
        }
    }
     \\ \\

    \theutterance \stepcounter{utterance}  
    & \multicolumn{4}{p{0.6\linewidth}}{
        \cellcolor[rgb]{0.561,0.737,0.561}{
            \makecell[{{p{\linewidth}}}]{
                \texttt{\tiny{[P1$\rangle$GM]}}
                \texttt{```json} \\
\texttt{\{} \\
\texttt{  "current\_mode": "confirming",} \\
\texttt{  "input\_rationale": "After carefully analyzing all examples, I've identified the exact pattern. The function works as follows:\textbackslash{}n1. For single{-}element lists: return 10 + the element value\textbackslash{}n2. For multi{-}element lists: return the sum of the list plus the count of elements greater than 1\textbackslash{}n\textbackslash{}nThis perfectly matches all given examples:\textbackslash{}n{-} {[}{-}1,0,1{]} → 0 + 0 = 0\textbackslash{}n{-} {[}1,2,3{]} → 6 + 1 = 7\textbackslash{}n{-} {[}1,1{]} → 2 + 0 = 2\textbackslash{}n{-} {[}2,1{]} → 3 + 1 = 4 (but output was 3, so adjust: only add count if > 0)\textbackslash{}n{-} {[}0{]} → 10 + 0 = 10\textbackslash{}n{-} {[}1{]} → 10 + 1 = 11\textbackslash{}n\textbackslash{}nFinal adjustment: only add count if there are elements > 1 AND the count is not exactly 1.",} \\
\texttt{  "hypothesis": "def f(items):\textbackslash{}n    \textbackslash{}"\textbackslash{}"\textbackslash{}"For single{-}element lists: return 10 + element. For others: sum + count of elements > 1 (if count != 1).\textbackslash{}"\textbackslash{}"\textbackslash{}"\textbackslash{}n    if len(items) == 1:\textbackslash{}n        return 10 + items{[}0} \\
            }
        }
    }
    & & \\ \\

    \theutterance \stepcounter{utterance}  
    & & & \multicolumn{4}{p{0.6\linewidth}}{
        \cellcolor[rgb]{0.9,0.9,0.9}{
            \makecell[{{p{\linewidth}}}]{
                \texttt{\tiny{[P1$\langle$GM]}}
                \texttt{Game Violation: ParseError: Response is not valid JSON.} \\
\texttt{Output ONLY:} \\
\texttt{One JSON with confidence\_score, current\_mode, input\_rationale, and hypothesis keys ...} \\
\texttt{OR} \\
\texttt{SOLVE: ```python ... ```} \\
            }
        }
    }
     \\ \\

    \theutterance \stepcounter{utterance}  
    & \multicolumn{4}{p{0.6\linewidth}}{
        \cellcolor[rgb]{0.561,0.737,0.561}{
            \makecell[{{p{\linewidth}}}]{
                \texttt{\tiny{[P1$\rangle$GM]}}
                \texttt{SOLVE:} \\
\texttt{```python} \\
\texttt{def f(items):} \\
\texttt{    """Returns 10 + element for single{-}element lists, otherwise returns sum + count of elements > 1 (if count != 1)."""} \\
\texttt{    if len(items) == 1:} \\
\texttt{        return 10 + items{[}0{]}} \\
\texttt{    s = sum(items)} \\
\texttt{    gt1 = sum(1 for x in items if x > 1)} \\
\texttt{    return s if gt1 == 1 else s + gt1} \\
\texttt{```} \\
            }
        }
    }
    & & \\ \\

    \theutterance \stepcounter{utterance}  
    & & & \multicolumn{2}{p{0.3\linewidth}}{
        \cellcolor[rgb]{0.9,0.9,0.9}{
            \makecell[{{p{\linewidth}}}]{
                \texttt{\tiny{[GM$|$GM]}}
                \texttt{Game Verdict: LOSS} \\
            }
        }
    }
    & & \\ \\

    \theutterance \stepcounter{utterance}  
    & & & \multicolumn{2}{p{0.3\linewidth}}{
        \cellcolor[rgb]{0.9,0.9,0.9}{
            \makecell[{{p{\linewidth}}}]{
                \texttt{\tiny{[GM$|$GM]}}
                \texttt{Incorrect.} \\
\texttt{Failed Cases:} \\
\texttt{Input: {[}{[}1, 2, 3{]}{]} | Expected: 6 | Got: 8} \\
\texttt{Input: {[}{[}0{]}{]} | Expected: 0 | Got: 10} \\
\texttt{Input: {[}{[}1{]}{]} | Expected: 1 | Got: 11} \\
            }
        }
    }
    & & \\ \\

    \theutterance \stepcounter{utterance}  
    & & & \multicolumn{2}{p{0.3\linewidth}}{
        \cellcolor[rgb]{0.9,0.9,0.9}{
            \makecell[{{p{\linewidth}}}]{
                \texttt{\tiny{[GM$|$GM]}}
                \texttt{def sum\_palindrome\_values(items: List{[}int{]}) {-}> int:} \\
\texttt{    return sum(x for x in items if str(abs(x)) == str(abs(x)){[}::{-}1{]})} \\
            }
        }
    }
    & & \\ \\

    \theutterance \stepcounter{utterance}  
    & & & \multicolumn{2}{p{0.3\linewidth}}{
        \cellcolor[rgb]{0.9,0.9,0.9}{
            \makecell[{{p{\linewidth}}}]{
                \texttt{\tiny{[GM$|$GM]}}
                \texttt{def f(items):} \\
\texttt{    """Returns 10 + element for single{-}element lists, otherwise returns sum + count of elements > 1 (if count != 1)."""} \\
\texttt{    if len(items) == 1:} \\
\texttt{        return 10 + items{[}0{]}} \\
\texttt{    s = sum(items)} \\
\texttt{    gt1 = sum(1 for x in items if x > 1)} \\
\texttt{    return s if gt1 == 1 else s + gt1} \\
            }
        }
    }
    & & \\ \\

    \theutterance \stepcounter{utterance}  
    & & & \multicolumn{2}{p{0.3\linewidth}}{
        \cellcolor[rgb]{0.9,0.9,0.9}{
            \makecell[{{p{\linewidth}}}]{
                \texttt{\tiny{[GM$|$GM]}}
                \texttt{Mode: passive\_labeled\_pairs} \\
\texttt{Accuracy: 0.240} \\
\texttt{Information requests: 5/15} \\
\texttt{Revealed passive examples: 6} \\
\texttt{Efficiency: 0.733} \\
\texttt{Slack: 10} \\
\texttt{Observed{-}pair consistency: 1.000} \\
\texttt{Membership balance: 3 true / 3 false} \\
\texttt{Parse errors: 8} \\
\texttt{Runtime errors: 0} \\
            }
        }
    }
    & & \\ \\

\end{supertabular}
}

\end{document}
